%! AP Statistics — Unit 4, Lesson 4.9 — Homework
\documentclass[10pt]{article}
\usepackage{apstats-article}
\usepackage{apstats-boxes}
\usepackage{grading}
\gradingassignment{apstats-u04-l09-hw}
\gradingcourse{advanced_statistics}
\newcommand{\TallMath}[1]{$\displaystyle #1\rule[-1.4em]{0pt}{3.2em}$}

\begin{document}

\pageheader{Unit 4, Lesson 4.9}{Homework}
{\small\itshape Write each final answer in its box. Only what is inside a box is graded.}\par

\vspace{0.1in}

\begin{practicebox}
\small

\textbf{1.} A random variable $X$ has $\mu_X = 20$ and $\sigma_X = 4$.
Define $W = 3 + 5X$.

\begin{enumerate}[label=(\alph*), leftmargin=2em, itemsep=8pt]
  \item $\mu_W = $ \makebox[4cm][l]{\hspace{1.6cm}\smash{\raisebox{-4pt}{\gradedblank[id=hw1a, type=short, points=1, answer=103, check=numeric, tolerance=0]{2.4cm}{0.2in}}}}
  \item $\sigma_W = $ \makebox[4cm][l]{\hspace{1.6cm}\smash{\raisebox{-4pt}{\gradedblank[id=hw1b, type=short, points=1, answer=20, check=numeric, tolerance=0]{2.4cm}{0.2in}}}}
  \item If $\mu_X$ increases by 2, what happens to $\mu_W$? \makebox[5cm][l]{\hspace{1.6cm}\smash{\raisebox{-4pt}{\gradedblank[id=hw1c, type=short, points=1, answer=10, check=numeric, tolerance=0, display={increases by 10}]{3.4cm}{0.2in}}}}
\end{enumerate}

\vspace{0.14in}

\textbf{2.} $X$ = cost of a large pizza (dollars). $Y$ = cost of a 2-liter soda (dollars).
The two items are purchased independently.
\begin{center}
\renewcommand{\arraystretch}{1.3}
\begin{tabular}{@{}lcc@{}}
& \textbf{Mean} & \textbf{SD} \\
\hline
$X$ & \$12.00 & \$2.00 \\
$Y$ & \$3.00  & \$0.75 \\
\end{tabular}
\end{center}

\begin{enumerate}[label=(\alph*), leftmargin=2em, itemsep=8pt]
  \item Find $\mu_{X+Y}$ and interpret in context.\hfill\gradedright{\smash{\raisebox{2pt}{$\mu_{X+Y} = $~\gradedblank[id=hw2a1, type=short, points=1, answer=15, check=numeric, tolerance=0.01, display={\$15.00}]{0.9in}{0.2in}}}}\\[4pt]
        \gradedhide{\gradedhere{\gradedblank[id=hw2a2, type=rubric, points=1, display={On average, a customer spends about \$15.00 on a pizza and soda combined.}]{\linewidth+6pt}{\gradedlinesheight{10}}}\writelines{10}}

  \item Find $\sigma_{X+Y}$. Show all work (add variances first!).\gradedbelow[r]{0.3in-0.25in}{$\sigma_{X+Y} = $~\gradedblank[id=hw2b, type=short, points=2, answer=2.14, check=numeric, tolerance=0.01]{0.9in}{0.2in}}

        \vspace{0.3in}

  \item A customer's friend claims the standard deviation of the total is $\$2.75$.
        Explain the error.\\
        \gradedhide{\gradedhere{\gradedblank[id=hw2c, type=rubric, points=2, display={SDs don't add: add variances, then take the square root.}]{\linewidth+6pt}{\gradedlinesheight{10}}}\writelines{10}}
\end{enumerate}

\vspace{0.14in}

\textbf{3.} Two independent bowling scores: $A$ with $\mu_A = 140$, $\sigma_A = 18$;
$B$ with $\mu_B = 120$, $\sigma_B = 15$.

\begin{enumerate}[label=(\alph*), leftmargin=2em, itemsep=8pt]
  \item Find $\mu_{A-B}$ and $\sigma_{A-B}$.\gradedbelow[r]{0.3in-0.25in}{$\mu_{A-B} = $~\gradedblank[id=hw3a1, type=short, points=1, answer=20, check=numeric, tolerance=0]{0.9in}{0.2in}\quad $\sigma_{A-B} = $~\gradedblank[id=hw3a2, type=short, points=2, answer=23.4, check=numeric, tolerance=0.1]{0.9in}{0.2in}}

        \vspace{0.3in}

\tcbbreak
  \item Interpret $\mu_{A-B}$ in context.\\
        \gradedhide{\gradedhere{\kern0.4in\gradedblank[id=hw3b, type=rubric, points=2, display={On average, player A scores about 20 pins more than player B per game.}]{\linewidth+6pt-0.4in}{\gradedlinesheight{2}-5pt}}\writelines{10}}

  \item Why isn't $\sigma_{A-B} = \sigma_A - \sigma_B = 3$?\\
        \gradedhide{\gradedhere{\gradedblank[id=hw3c, type=rubric, points=2, display={Variances add even for a difference: both add variability, so $\sigma_{A-B} = 23.4$.}]{\linewidth+6pt}{\gradedlinesheight{10}}}\writelines{10}}
\end{enumerate}

\vspace{0.14in}

\textbf{4.} AP-Style. $X$ and $Y$ are independent. $\mu_X = 5$, $\sigma_X = 2$,
$\mu_Y = 3$, $\sigma_Y = 4$.

\begin{enumerate}[label=(\alph*), leftmargin=2em, itemsep=8pt]
  \item $\mu_{2X-Y} = $ \makebox[4cm][l]{\hspace{1.6cm}\smash{\raisebox{-4pt}{\gradedblank[id=hw4a, type=short, points=1, answer=7, check=numeric, tolerance=0]{2.4cm}{0.2in}}}}
  \item $\sigma_{2X-Y} = $ \makebox[4cm][l]{\hspace{1.6cm}\smash{\raisebox{-4pt}{\gradedblank[id=hw4b, type=short, points=2, answer=5.66, check=numeric, tolerance=0.01]{2.3cm}{0.2in}}}} (show work)

        \vspace{0.25in}
\end{enumerate}

\end{practicebox}

\vspace{0.1in}

\begin{extensionbox}
\small
\textbf{Extension --- optional.}
$X \sim N(100, 15)$ and $Y \sim N(80, 12)$ are independent (scores on two different tests).
\begin{enumerate}[leftmargin=1.5em, itemsep=6pt]
  \item Find $\mu_{X+Y}$ and $\sigma_{X+Y}$. State the distribution of $T = X+Y$.

        \vspace{0.25in}

  \item Find $P(T > 200)$. Show work.

        \vspace{0.35in}
\end{enumerate}
\end{extensionbox}

\vspace{0.1in}

\begin{blurbbox}[Preview: Lesson 4.10]{sky}
\small
In Lesson 4.10 you will meet the \textbf{binomial distribution} --- a special discrete
distribution for the count of ``successes'' in $n$ independent trials. Its mean $\mu = np$
and variance $\sigma^2 = np(1-p)$ follow directly from the rules you learned today.
\end{blurbbox}

\end{document}
